\documentclass[10pt,a4paper]{amsart}
\usepackage[margin=19mm]{geometry}
\usepackage{amsmath,amssymb,mathtools}
\usepackage[scr=rsfs,cal=euler]{mathalfa}
\usepackage{xfrac}
\usepackage[unicode,hidelinks,bookmarks=false]{hyperref}
\newcommand{\ie}{i.e.\ }
\newcommand{\cf}{cf.\ }
\newcommand{\resp}{respectively\ }
\newcommand{\Subsection}{\subsection}
\providecommand{\nobreakdash}{\nobreak\hbox{-}}
\setlength{\emergencystretch}{2em}
\multlinegap0pt
\usepackage{amssymb}
\usepackage{xcolor}
\newtheorem{cor}{}
\newtheorem{lem}{}
\newcommand{\M}{\mathcal{M}}
\newcommand{\R}{\mathbb{R}} 
\title{Eyring-Kramers Law for the Underdamped Langevin Process}
\author{Seungwoo Lee, Mouad Ramil, Insuk Seo}
\date{Source: arXiv:2503.12610v2 (CC BY 4.0)}
\begin{document}
\maketitle

\noindent\emph{Source and licence.} Eyring-Kramers Law for the Underdamped Langevin Process. Version arXiv:2503.12610v2 (CC BY 4.0), distributed by arXiv under the Creative Commons Attribution 4.0 International licence, http://creativecommons.org/licenses/by/4.0/, as recorded in the arXiv licence metadata for this exact version and shown on the abstract page https://arxiv.org/abs/2503.12610v2. Full licence notice: \texttt{licenses/arxiv-2503.12610-CC-BY-4.0.txt}.\par\smallskip

\noindent\textit{Editorial scope.} Counted unit: the article's lem velocity (source0.tex lines 1455--1461). The article's complete original proof of it is delivered with it (source0.tex lines 1462--1537, one \texttt{proof} environment of 76 lines). No mathematics is written, repaired or re-derived: the statement and the proof are the article's own lines, in the article's own notation, and no retained line is substituted or re-flowed.  Retained uncounted as the source-local context the proof needs: lines 201--212; lines 853--856; lines 1423--1426.  Nothing was omitted from any retained span, so the package carries no omission record.  No retained line contains a citation, so the wrapper carries no bibliography and no bibliography entry.  No retained line contains a figure environment, an image-inclusion command or drawing code, so no image is delivered and the package declares no asset dependency.  Editorial formatting only, in addition: an AMS \texttt{amsart} preamble that restates the article's own declarations and macros used by the retained lines, namely the environments \texttt{cor}, \texttt{lem} on separate counters, together with the standard packages those lines need.  The retained environments print in this document as Corollary 1, Lemma 1 in order of appearance, whereas the article prints the same items with the numbering of its own full document.  The complete native source of the article as distributed in the arXiv e-print for this exact version is bundled unchanged as \texttt{sources/174868/source0.tex}.
To describe this model concretely, we repeat the definition here. Let $d\geq1$, and we shall focus on the study of the underdamped Langevin process $(X^{\epsilon}(t)=(q^{\epsilon}(t),\,p^{\epsilon}(t)))_{t\geq0}$
in $\mathbb{R}^{d}\times\mathbb{R}^{d}$ which is solution to the following SDE: 
\begin{equation}
\left\{ \begin{aligned} & \mathrm{d}q^{\epsilon}(t)=p^{\epsilon}(t)\mathrm{d}t\;,\\
 & \mathrm{d}p^{\epsilon}(t)=-\nabla U(q^{\epsilon}(t))\mathrm{d}t-\gamma p^{\epsilon}(t)\mathrm{d}t+\sqrt{2\gamma\epsilon}\mathrm{d}B(t)\;,
\end{aligned}
\right.\label{eq:sde}
\end{equation}
where $U\in C^{2}(\mathbb{R}^{d},\mathbb{\,R})$ is the potential function,
$\gamma>0$ and $\epsilon>0$ denote the friction coefficient and
the temperature, respectively. In order to prove the Eyring-Kramers law for this process, we will investigate its behaviour when the temperature
$\epsilon>0$ is small.

\begin{align}
\mathcal{O}_{N,\,\epsilon} & :=\mathcal{S}_{\epsilon}\cup\{x\in\mathbb{R}^{2d}:|x|>N\}\;\;\text{and}\label{eq:definition O_N}\\
\mathcal{D}_{N,\,\epsilon} & :=\textup{B}(0,\,N)\setminus(\overline{\mathcal{M}_{\epsilon}}\cup\overline{\mathcal{S}_{\epsilon}})\;.\nonumber 
\end{align}

\begin{cor}\label{cor:weyl density}
There exists a constant $C>0$ such that for all $t\in(0,r]$, if $\eta,\,\eta',\,\epsilon>0$ are small enough, then
$$\sup_{x\in\partial\M_\epsilon}\mathbb{P}_{x}\left(\big|q^{\epsilon}(t)-s\big|\in(\epsilon-\eta,\epsilon+\eta),\,\big|p^{\epsilon}(t)\big|\leq\eta',\,\tau_{\mathcal{O}_{N,\,\epsilon}}^{\epsilon}>t\right)\leq C\eta\,\eta'.$$
\end{cor}

\begin{lem}
\label{lem:control inward velocity} For $\epsilon$ small enough
and $N$ large enough, 
\[
\sup_{x\in\partial\mathcal{M}_{\epsilon}}\mathbb{P}_{x}\left(|p^{\epsilon}(\tau_{\mathcal{O}_{N,\,\epsilon}}^{\epsilon})|\leq\delta,\tau_{\mathcal{O}_{N,\,\epsilon}}^{\epsilon}\leq r\right)\underset{\delta\rightarrow0}{\longrightarrow}0\;.
\]
\end{lem}

\begin{proof}

Let $\delta\in(0,1)$. Given the definition of $\mathcal{O}_{N,\,\epsilon}$ in~\eqref{eq:definition O_N}, either $X^{\epsilon}(\tau_{\mathcal{O}_{N,\,\epsilon}}^{\epsilon})\in\partial\mathcal{S}_{\epsilon}$ or $X^{\epsilon}(\tau_{\mathcal{O}_{N,\,\epsilon}}^{\epsilon})\in\partial\mathrm{B}(0,N)$. Therefore, let us start by considering the first case and we prove that for any $\nu>0$, if $\delta$ is small enough then 
\begin{equation}
\sup_{x\in\partial\mathcal{M}_{\epsilon}}\mathbb{P}_{x}\left(X^{\epsilon}(\tau_{\mathcal{O}_{N,\,\epsilon}}^{\epsilon})\in\partial\mathcal{S}_{\epsilon},\,|p^{\epsilon}(\tau_{\mathcal{O}_{N,\,\epsilon}}^{\epsilon})|\leq\delta,\,\tau_{\mathcal{O}_{N,\,\epsilon}}^{\epsilon}\leq r\right)\leq\nu\;.\label{eq:sup proba inward velocity proof}
\end{equation}
Let $\eta\in(0,\,1/2)$ and let
\[
\mathcal{Z}:=\sup_{0\leq t<t'\leq r}\frac{|B(t)-B(t')|}{|t-t'|^{\eta}},
\]
then $\mathcal{Z}<\infty$, almost-surely, since the Brownian motion
$(B(t))_{t\geq0}$ is $\eta$-H\"olderian.

Let $J_\delta=\lfloor1/\delta\rfloor$ and let $r_{k}=r\frac{k}{J_\delta}$ for $0\leq k\leq J_\delta$. We take $M>0$ large enough such that $\mathbb{P}(\mathcal{Z}>M)\leq\nu/2$, then 
\begin{align}
 & \mathbb{P}_{x}\left(X^{\epsilon}(\tau_{\mathcal{O}_{N,\,\epsilon}}^{\epsilon})\in\partial\mathcal{S}_{\epsilon},\,|p^{\epsilon}(\tau_{\mathcal{O}_{N,\,\epsilon}}^{\epsilon})|\leq\delta,\,\tau_{\mathcal{O}_{N,\,\epsilon}}^{\epsilon}\leq r\right)\nonumber \\
 & \leq\mathbb{P}_{x}\left(X^{\epsilon}(\tau_{\mathcal{O}_{N,\,\epsilon}}^{\epsilon})\in\partial\mathcal{S}_{\epsilon},\,|p^{\epsilon}(\tau_{\mathcal{O}_{N,\,\epsilon}}^{\epsilon})|\leq\delta,\,\tau_{\mathcal{O}_{N,\,\epsilon}}^{\epsilon}\leq r,\,\mathcal{Z}\leq M\right)+\mathbb{P}(\mathcal{Z}>M)\nonumber \\
 & \leq\sum_{k=0}^{J_\delta-1}\mathbb{P}_{x}\left(\tau_{\mathcal{O}_{N,\,\epsilon}}^{\epsilon}\in(r_{k},\,r_{k+1}],\,X^{\epsilon}(\tau_{\mathcal{O}_{N,\,\epsilon}}^{\epsilon})\in\partial\mathcal{S}_{\epsilon},\,|p^{\epsilon}(\tau_{\mathcal{O}_{N,\,\epsilon}}^{\epsilon})|\leq\delta,\,\mathcal{Z}\leq M\right)+\frac{\nu}{2}\;.\label{eq:sum interval exit time}
\end{align}

Moreover, by~\eqref{eq:sde}, for all $u\in[r_{k},\,\tau_{\mathcal{O}_{N,\,\epsilon}}^{\epsilon}]$,
\[
p^{\epsilon}(\tau_{\mathcal{O}_{N,\,\epsilon}}^{\epsilon})-p^{\epsilon}(u)=\int_{u}^{\tau_{\mathcal{O}_{N,\,\epsilon}}^{\epsilon}}\left[-\nabla U(q^{\epsilon}(\rho))-\gamma p^{\epsilon}(\rho)\right]\mathrm{d}\rho+\sqrt{2\gamma\epsilon}\left(B(\tau_{\mathcal{O}_{N,\,\epsilon}}^{\epsilon})-B(u)\right)\;.
\]
Under the event $\{\tau_{\mathcal{O}_{N,\,\epsilon}}^{\epsilon}\in(r_{k},\,r_{k+1}],\,\mathcal{Z}\leq M\}$, using the boundedness of $\R^{2d}\setminus\mathcal{O}_{N,\,\epsilon}$ there exists a constant $C=C(N)\geq1$ such that
\begin{equation}\label{eq:triangle ineq velocity boundary}
\big|p^{\epsilon}(\tau_{\mathcal{O}_{N,\,\epsilon}}^{\epsilon})-p^{\epsilon}(u)\big|\leq\frac{C}{J_\delta}+\frac{C}{J_\delta^\eta}\leq \frac{2C}{J_\delta^\eta}\;.
\end{equation}
In particular, under the event $\{\tau_{\mathcal{O}_{N,\,\epsilon}}^{\epsilon}\in(r_{k},\,r_{k+1}],\,|p^{\epsilon}(\tau_{\mathcal{O}_{N,\,\epsilon}}^{\epsilon})|\leq\delta,\,\mathcal{Z}\leq M\}$, one has by the triangle inequality,
\begin{align}
\big|p^{\epsilon}(r_k)\big|&\leq\big|p^{\epsilon}(\tau_{\mathcal{O}_{N,\,\epsilon}}^{\epsilon})-p^{\epsilon}(r_k)\big|+\big|p^{\epsilon}(\tau_{\mathcal{O}_{N,\,\epsilon}}^{\epsilon})\big|\nonumber\\
&\leq\frac{2C}{J_\delta^\eta}+\frac{1}{J_\delta}\leq\frac{3C}{J_\delta^\eta}\label{eq:ineq triangl p}\,,
\end{align}
since $\delta\leq1/J_\delta$, by definition of $J_\delta$. Moreover, under the event $\{\tau_{\mathcal{O}_{N,\,\epsilon}}^{\epsilon}\in(r_{k},\,r_{k+1}],\,|p^{\epsilon}(\tau_{\mathcal{O}_{N,\,\epsilon}}^{\epsilon})|\leq\delta,\,\mathcal{Z}\leq M\}$, the triangle inequality also yields by~\eqref{eq:triangle ineq velocity boundary},
\begin{align}
\big|q^{\epsilon}(\tau_{\mathcal{O}_{N,\,\epsilon}}^{\epsilon})-q^{\epsilon}(r_{k})\big| & =\left|\int_{r_{k}}^{\tau_{\mathcal{O}_{N,\,\epsilon}}^{\epsilon}}p^{\epsilon}(u)\mathrm{d}u\right|\nonumber\\
&\leq \frac{2C}{J_\delta^{1+\eta}}+\int_{r_{k}}^{\tau_{\mathcal{O}_{N,\,\epsilon}}^{\epsilon}}\big|p^{\epsilon}(\tau_{\mathcal{O}_{N,\,\epsilon}}^{\epsilon})\big|\mathrm{d}u\nonumber\\ 
 &\leq\frac{2C}{J_\delta^{1+\eta}}+\frac{1}{J_\delta^{2}}\leq\frac{3C}{J_\delta^{1+\eta}}\label{eq:first triangle ineq}\,.
\end{align}
Furthermore, on the event $\{X^{\epsilon}(\tau_{\mathcal{O}_{N,\,\epsilon}}^{\epsilon})\in\partial\mathcal{S}_{\epsilon}\}$, 
\begin{equation}\label{eq:equation boundary S}
\big|q^{\epsilon}(\tau_{\mathcal{O}_{N,\,\epsilon}}^{\epsilon})-s\big|^{2}+\big|p^{\epsilon}(\tau_{\mathcal{O}_{N,\,\epsilon}}^{\epsilon})\big|^{2}=\epsilon^2\;.
\end{equation}
Therefore, on the event $\{X^{\epsilon}(\tau_{\mathcal{O}_{N,\,\epsilon}}^{\epsilon})\in\partial\mathcal{S}_{\epsilon},\,\big|p^{\epsilon}(\tau_{\mathcal{O}_{N,\,\epsilon}}^{\epsilon})\big|\leq\delta\}$,
\[
\sqrt{\epsilon^2-\frac{1}{J^2_\delta}}\leq\big|q^{\epsilon}(\tau_{\mathcal{O}_{N,\,\epsilon}}^{\epsilon})-s\big|\leq \epsilon\;.
\]
Combining the inequality above with~\eqref{eq:first triangle ineq} one obtains
\[
\sqrt{\epsilon^2-\frac{1}{J^2_\delta}}-\frac{3C}{J_\delta^{1+\eta}}\leq\big|q^{\epsilon}(r_k)-s\big|\leq \epsilon+\frac{3C}{J_\delta^{1+\eta}}\;.
\]
Taking $\delta$ sufficiently small and applying the asymptotic expansion of $\sqrt{1-x}$ when $x\rightarrow0$, one obtains the existence of a constant $C=C(N,\epsilon)>0$ which is independent of $\delta$, such that
\[
\epsilon-\frac{C}{J_\delta^{1+\eta}}\leq\big|q^{\epsilon}(r_k)-s\big|\leq \epsilon+\frac{C}{J_\delta^{1+\eta}}\;.
\]

All in all, combining the inequality above with~\eqref{eq:ineq triangl p} and taking a constant $C>0$ large enough, one obtains for all $0\leq k\leq J_\delta-1$ and $x\in\partial\M_\epsilon$,
\begin{align*}
&\mathbb{P}_{x}\left(\tau_{\mathcal{O}_{N,\,\epsilon}}^{\epsilon}\in(r_{k},\,r_{k+1}],\,X^{\epsilon}(\tau_{\mathcal{O}_{N,\,\epsilon}}^{\epsilon})\in\partial\mathcal{S}_{\epsilon},\,|p^{\epsilon}(\tau_{\mathcal{O}_{N,\,\epsilon}}^{\epsilon})|\leq\delta,\,\mathcal{Z}\leq M\right)\nonumber \\
 & \leq\mathbb{P}_{x}\left(\tau_{\mathcal{O}_{N,\,\epsilon}}^{\epsilon}\in(r_{k},\,r_{k+1}],\,\big|q^{\epsilon}(r_{k})-s\big|\in\big[\epsilon-\frac{C}{J_\delta^{1+\eta}},\,\epsilon+\frac{C}{J_\delta^{1+\eta}}\big],\,\big|p^{\epsilon}(r_{k})\big|\leq\frac{C}{J^{\eta}_\delta}\right)\;.
\end{align*}
In particular, for $k=0$, since the initial condition $x\in\partial\M_\epsilon$, it is easy to see that, taking $\epsilon$ and $\delta$ sufficiently small, ensures that the probability above vanishes. Therefore, it is sufficient to consider the case $1\leq k\leq J_\delta-1$. Moreover, by Corollary~\ref{cor:weyl density}, there exists a constant $C'>0$ independent of $k,\delta$ and $x\in\partial\M_\epsilon$ such that
\begin{align*}
\mathbb{P}_{x}\left(\tau_{\mathcal{O}_{N,\,\epsilon}}^{\epsilon}\in(r_{k},\,r_{k+1}],\,\big|q^{\epsilon}(r_{k})-s\big|\in\big[\epsilon-\frac{C}{J_\delta^{1+\eta}},\,\epsilon+\frac{C}{J_\delta^{1+\eta}}\big],\,\big|p^{\epsilon}(r_{k})\big|\leq\frac{C}{J^{\eta}_\delta}\right) 
 \leq \frac{C'}{J_\delta^{1+2\eta}}\,.
\end{align*}
Reinjecting in~\eqref{eq:sum interval exit time} and summing over $1\leq k\leq J_\delta-1$, it follows that
\begin{align}
 \mathbb{P}_{x}\left(X^{\epsilon}(\tau_{\mathcal{O}_{N,\,\epsilon}}^{\epsilon})\in\partial\mathcal{S}_{\epsilon},\,|p^{\epsilon}(\tau_{\mathcal{O}_{N,\,\epsilon}}^{\epsilon})|\leq\delta,\,\tau_{\mathcal{O}_{N,\,\epsilon}}^{\epsilon}\leq r\right)\nonumber \leq \frac{C'}{J_\delta^{2\eta}}+\frac{\nu}{2}\;.
\end{align}
Taking $\delta$ sufficiently small concludes the proof. It remains however to also show that 
\[
\sup_{x\in\partial\mathcal{M}_{\epsilon}}\mathbb{P}_{x}\left(X^{\epsilon}(\tau_{\mathcal{O}_{N,\,\epsilon}}^{\epsilon})\in\partial\mathrm{B}(0,\,N),\,|p^{\epsilon}(\tau_{\mathcal{O}_{N,\,\epsilon}}^{\epsilon})|\leq\delta,\,\tau_{\mathcal{O}_{N,\,\epsilon}}^{\epsilon}\leq r\right)\underset{\delta\rightarrow0}{\longrightarrow}0\;.
\]
Nonetheless, a similar proof also applies in this case by replacing the equality~\eqref{eq:equation boundary S} by the appropriate formula on $\partial\mathrm{B}(0,\,N)$, hence the proof of this lemma. 
\end{proof}

\end{document}
